% Purpose: PhD Dissertation Abstract summarizing research motivations, methods, benchmark datasets, HPC implementation, and scientific outcomes.
% Status: Draft; under academic review.
% Source-of-truth: OGS/src/, doc/MHPC/chapters/abstract.tex, experimental results.

\begin{abstract}
The exponential expansion of continuous seismic recording archives and the deployment of dense regional networks have created severe operational bottlenecks for traditional earthquake monitoring workflows. Classical automated phase pickers, such as Short-Time Average over Long-Time Average (\ac{STA/LTA}) and autoregressive Akaike Information Criterion (\ac{AIC}) formulations, degrade substantially when confronted with low signal-to-noise ratios, emergent phase onsets, overlapping seismic codas, and anthropogenic noise. Consequently, regional seismic observatories remain heavily dependent on labor-intensive manual analyst review, introducing latency and limiting catalog completeness at microseismic scales.

This dissertation presents the design, implementation, and empirical validation of an end-to-end, procedural machine learning pipeline for automated earthquake pick detection, phase association, hypocentral relocation, and magnitude characterization, specifically engineered for scalable execution on modern \ac{HPC} architectures. The computational framework leverages the open-source SeisBench ecosystem to deploy deep-learning phase picking architectures, including the 1D U-Net convolutional network PhaseNet and the hierarchical convolutional-transformer network \ac{EQTransformer}. To systematically investigate domain adaptation and model transferability, these pickers are evaluated across diverse regional and global benchmark training corpora: the Italian seismic dataset \ac{INSTANCE}, the global \ac{STEAD} corpus, the Southern California \ac{SCEDC} archive, and the original author training sets.

The target case study comprises continuous three-component waveform streams recorded by the \acf{SMINO}, managed by the \acf{OGS}, augmented by seventeen regional and transboundary networks in Veneto, Trentino-Alto Adige, Austria, and Slovenia. The network configuration encompasses 266 stations and over 790 channels sampled at 100~Hz. The observation horizon spans a 93-day window from March 20 to June 20, 2024, centered on the March 27, 2024, $M_{\mathrm{L}}~4.6$ ($M_{\mathrm{w}}~4.5$) Socchieve earthquake in Carnia, Friuli—one of the most energetic seismic episodes in Northeastern Italy since the destructive 1976 sequence.

Detected phase arrivals are clustered and associated into discrete seismic events using two complementary state-of-the-art algorithms: the Bayesian Gaussian Mixture Model Associator (\ac{GaMMA}) and the high-throughput octree-based associator \ac{PyOcto}. Associated events are relocated using a 1D probabilistic non-linear search via \ac{NonLinLoc}, followed by parametric local magnitude ($M_{\mathrm{L}}$) inversion. To establish an unbiased comparative baseline, the OGS reference catalog for the observation window was independently relocated under the identical 1D \ac{NonLinLoc} velocity model, yielding 380 confirmed events and 7,129 manually curated phase picks (4,126~P and 3,003~S). Performance is evaluated using a physically bounded bipartite graph maximum-weight matching framework (\ac{BPGMA}), resolving matched, missed, and proposed events and picks with strict spatiotemporal and phase-type conservation.

Distributed execution on the CINECA Leonardo supercomputer (EuroHPC Booster partition, utilizing NVIDIA A100 GPUs and multi-threaded CPU pipelines) demonstrates near-linear scalability, reducing continuous waveform processing times by orders of magnitude compared to legacy serial workflows. Systematic evaluation reveals that PhaseNet trained on the regional \ac{INSTANCE} dataset achieves the highest pick recovery, optimal phase residual distributions, and the most reliable association completeness, demonstrating that domain-specific regional noise representation outweighs sheer dataset volume for regional seismotectonic surveillance. The proposed infrastructure provides operational observatories with a scalable, reproducible foundation for both real-time early warning alerting and high-resolution historical catalog reprocessing.
\end{abstract}
